\section{Two double covers and four local counts}\label{sec:intro}

Take two independent unramified double covers of a smooth curve $C$.
Each can acquire a normal deformation while $C$ itself stays rigid.
Suppose the first deformation occurs on one axis of a parameter plane
and the second on the other. Crossing the first axis changes the local
curve count. How much does that change depend on which side of the
second axis we cross? The answer is a local invariant on the connected
fourfold cover formed by taking the fiber product of the two doubles.

We hold the curve $C$ and its complex rank-two normal bundle $N$ fixed,
where $C$ has genus $g\ge2$ and $\deg N=2g-2$. The parameter varies a
real Cauchy--Riemann operator $D$ on $N$. A section in the kernel of
its pullback to a cover is an infinitesimal normal deformation there.
At a parameter where the pullback operator has zero kernel on every
connected stable map of positive degree at most $M$, the
zero section has a local Gromov--Witten contribution in each degree
$d\le M$ and each arithmetic source genus $h$. Write it as $L^C_{d,h}$.
It counts connected stable maps to $C$ with the Euler class of their
normal obstruction bundle. The covers on which we prescribe a kernel
are unramified. The stable maps in the count can have branch points,
nodes, and contracted components.

Let $Z$ be the common fourfold cover. Theorem~\ref{thm:square}
constructs a family with the two prescribed walls. Label its chambers
by the signs of the two parameters. For $d\le M$, in a sufficiently
small rectangle about the crossing, one has
\[
  L^C_{d,h}(++)-L^C_{d,h}(-+)-L^C_{d,h}(+-)+L^C_{d,h}(--)
                    = L^Z_{d/4,h}(++) .
\]
The right side is zero unless four divides $d$. Its operator is the
pullback of the positive-corner operator. A map changes sign across
an axis precisely when it lifts to the corresponding double. In the
alternating difference only maps lifting to both doubles survive,
each with multiplier four. Such a map has exactly four lifts to $Z$
with its domain fixed. This accounts for the coefficient one on the
right, including the automorphism weights of the maps. For example,
over a curve of genus two the first common lift has degree four and
genus five. Its contribution to the difference in that genus is its
determinant sign.

\subsection*{From signs on covers to an equality of invariants}
Parker computed the determinant signs carried by the irreducible real
representations of a covering group \cite[\S1.2]{parker}.
Bryan and Pandharipande stratified the space of maps to a curve by
the intermediate covers to which they lift \cite[\S2.2]{bp01}.
Wendl's super-rigidity theorem and his analysis of nonabelian
bifurcations, and Doan and Walpuski's study of the double-cover wall,
describe how normal kernels arise \cite[\S2.4]{wendl},
\cite[Theorem~1.3]{dw}. To obtain the displayed equality one must
realize both walls without other kernels intervening, then compare
the two normal Euler classes over the entire compactified space of maps.

There is a compactness issue even for the two doubles. A bound on the
degree does not bound the source genus, so a separate neighborhood for
each moduli space would give less than the formula asserts.
A normal section has a graph whose area is bounded by its covering
degree. A sequence of unexpected sections therefore has a limit as a
cycle. The sheets of its support are separated. Lifting to the nearest
sheet and averaging the sections there produces a kernel on a fixed
intermediate cover. The prescribed kernels account for all such
limits (Section~\ref{sec:family}). For the Euler class comparison we
fix the monodromy image of each map. Parallel kernel vectors match
across a node exactly when their values on its two branches agree.
Removing these vectors leaves a continuous reduced obstruction bundle
across the wall. Its orientation comparison supplies the determinant
sign in the integral (Section~\ref{sec:oriented}).

The same argument applies on a finite principal unramified $G$-cover.
The kernels may be prescribed by pairwise inequivalent nontrivial
irreducible real representations of real type, provided their
associated flat bundles together embed in the underlying real bundle
of $N$. On a map with monodromy image $H$, the sign of each crossing
is given by the parity of the dimension of its $H$-fixed subspace.
Writing these signs in the basis of transitive $G$-sets expresses the
mixed difference as a sum over intermediate covers
(Theorem~\ref{thm:mixed}). Section~\ref{sec:factors} proves the
corresponding identity for the local factors in the universal cycle
algebra, using the integral generators of \cite[Theorem~1.1]{WueGenerators}.
In the two-double example the mixed connected count in its least
source genus has exact denominator $p^j$ in degree $4p^j$, for every
prime $p$ and every $j\ge1$ (Theorem~\ref{thm:prime}).

\subsection*{Endpoints and examples}
At the minimum source genus the Riemann--Hurwitz formula leaves only
unramified maps. At two arbitrary rigid endpoints, the determinant signs at both ends
therefore give the change in that genus. Equivalently, multiplying
each initial covering sign by its endpoint ratio gives its final sign. In higher genus we prove a transport formula along
paths made up of the realized unramified crossings and rigid segments.
Its initial data are the full series on the intermediate covers
(Theorem~\ref{thm:endpoints}). Thus a path calculation can start
with an independently computed chamber. For positive maximal normal holonomy, Appendices~\ref{sec:maximal}--\ref{sec:flat-sign}
identify the full series when the orthogonal bundle is orientable and
compute the finite-cover signs in both cases. The absolute all-genus
calculation for a nonorientable orthogonal bundle is given in
\cite[Theorem~1.1]{WueSpin}.

For an elliptic base the suspension of a symplectic return map gives
a separate description. Its horizontal tori are the finite orbits of
the return map. When the finite orbits are isolated as a set in the
sense of Section~\ref{sec:elliptic-series}, the fixed-point indices
of the iterates recover the local cycle exponents by M\"obius inversion.
A projective example has sixteen torsion-distinct sections with the
same double class. In higher genus one generic path between specified
spin operators forces a bifurcation in every covering degree $m\ge2$
(Theorem~\ref{thm:walls}). Along a subsequence, the domains converge
to an infinite cover carrying a nonzero bounded solution of the
limiting operator (Proposition~\ref{prop:wall-limit}).
